% Shared macros for the urcu-txn paper series.
% Included by each paper's main.tex via % Shared macros for the urcu-txn paper series.
% Included by each paper's main.tex via % Shared macros for the urcu-txn paper series.
% Included by each paper's main.tex via % Shared macros for the urcu-txn paper series.
% Included by each paper's main.tex via \input{../common/macros}.
%
% Keep this class-agnostic: nothing here may assume `article`, so that a paper
% can be retargeted to acmart/IEEEtran by changing only its own preamble.

% ---------------------------------------------------------------------------
% Terminology.
%
% These exist so the series says the same thing the same way every time, and so
% that a terminology decision is changed in one place. The guarantee wording is
% load-bearing: this engine is NOT STM and is NOT serializable. See
% \pseudotxn below and the note in README.md.
% ---------------------------------------------------------------------------
\newcommand{\pseudotxn}{pseudo-transaction\xspace}
\newcommand{\pseudotxns}{pseudo-transactions\xspace}
\newcommand{\mcas}{MCAS\xspace}
\newcommand{\kcas}{$k$-CAS\xspace}
\newcommand{\rcu}{RCU\xspace}
\newcommand{\QSBR}{QSBR\xspace}
\newcommand{\urcu}{userspace RCU\xspace}
\newcommand{\sw}{single-writer\xspace}
\newcommand{\mw}{multi-writer\xspace}
\newcommand{\fliplatch}{flip-latch\xspace}
\newcommand{\ryw}{read-your-own-writes\xspace}

% Engine identifiers.
\newcommand{\urcutxn}{\texttt{urcu-txn}\xspace}
\newcommand{\rcumcas}{\texttt{rcu-mcas}\xspace}
\newcommand{\rcutxnsw}{\texttt{rcu-txn-sw}\xspace}

% ---------------------------------------------------------------------------
% Code and identifiers.
% ---------------------------------------------------------------------------
% \code offers TeX no place to break, and this series is dense with identifiers
% long enough to overflow a line on their own (cds_list_for_each_entry_reverse,
% -DURCU_TXN_SW_EXCL_VALIDATE). Two of them in one sentence is enough to push
% text into the margin no matter how the surrounding prose is stretched.
%
% Since every such identifier already writes its underscores as \_, redefining
% \_ locally gives TeX a legal breakpoint after each one -- silently, with NO
% hyphen. That last part is the whole reason not to reach for hyphenat's [htt]
% option instead: a hyphen inside a C identifier is ambiguous, and the reader
% cannot tell whether it belongs to the name. For a paper whose purpose is
% enablement, an identifier someone has to guess at is a real defect; a slightly
% loose line is not.
%
% Scoped to \texttt's own group, so \_ keeps its normal meaning everywhere else.
%
% \let, not \renewcommand, and \DeclareRobustCommand, not \newcommand: \code is
% used inside \caption, which is a MOVING argument (written to the .aux and
% re-read). A fragile \code expands there and the \renewcommand form dies with
% "Argument of \caption@ydblarg has an extra }" -- a fatal error whose message
% names neither \code nor the caption it came from.
\newcommand{\codeunderscore}{\textunderscore\allowbreak}
\DeclareRobustCommand{\code}[1]{\texttt{\let\_\codeunderscore #1}}
\DeclareRobustCommand{\fn}[1]{\code{#1()}}

% A record in a transaction's write set: {slot, old, new}.
\newcommand{\record}[3]{$\{#1,\,#2 \rightarrow #3\}$}
% An edit on one slot, old to new.
\newcommand{\edit}[2]{$#1 \rightarrow #2$}

% Transaction status values.
\newcommand{\UNDECIDED}{\textsc{undecided}\xspace}
\newcommand{\SUCCEEDED}{\textsc{succeeded}\xspace}
\newcommand{\FAILED}{\textsc{failed}\xspace}
\newcommand{\ABORT}{\textsc{abort}\xspace}

% ---------------------------------------------------------------------------
% Editorial markers. Define \draftmode in the preamble to make these visible;
% they vanish in a submission build. Never ship with \draftmode set.
% ---------------------------------------------------------------------------
\newif\ifdraftmode
\draftmodefalse

% These use \color/\bfseries SWITCHES inside a group rather than
% \textcolor{}{}/\textbf{}, whose arguments are short: an editorial note
% spanning a blank line would otherwise die with "Paragraph ended before
% \@textcolor was complete". Notes routinely span paragraphs, so the switch
% form is the only one that survives real use.
\newcommand{\marker}[3]{\ifdraftmode{\color{#1}\bfseries[#2: #3]}\fi}

\newcommand{\TODO}[1]{\marker{red}{TODO}{#1}}
\newcommand{\NOTE}[1]{\marker{blue}{NOTE}{#1}}
% Claims needing a citation before submission.
\newcommand{\NEEDCITE}[1]{\marker{orange}{CITE}{#1}}
% Numbers currently sourced from source comments rather than a reproducible run.
\newcommand{\UNVERIFIED}[1]{\marker{purple}{UNVERIFIED}{#1}}
.
%
% Keep this class-agnostic: nothing here may assume `article`, so that a paper
% can be retargeted to acmart/IEEEtran by changing only its own preamble.

% ---------------------------------------------------------------------------
% Terminology.
%
% These exist so the series says the same thing the same way every time, and so
% that a terminology decision is changed in one place. The guarantee wording is
% load-bearing: this engine is NOT STM and is NOT serializable. See
% \pseudotxn below and the note in README.md.
% ---------------------------------------------------------------------------
\newcommand{\pseudotxn}{pseudo-transaction\xspace}
\newcommand{\pseudotxns}{pseudo-transactions\xspace}
\newcommand{\mcas}{MCAS\xspace}
\newcommand{\kcas}{$k$-CAS\xspace}
\newcommand{\rcu}{RCU\xspace}
\newcommand{\QSBR}{QSBR\xspace}
\newcommand{\urcu}{userspace RCU\xspace}
\newcommand{\sw}{single-writer\xspace}
\newcommand{\mw}{multi-writer\xspace}
\newcommand{\fliplatch}{flip-latch\xspace}
\newcommand{\ryw}{read-your-own-writes\xspace}

% Engine identifiers.
\newcommand{\urcutxn}{\texttt{urcu-txn}\xspace}
\newcommand{\rcumcas}{\texttt{rcu-mcas}\xspace}
\newcommand{\rcutxnsw}{\texttt{rcu-txn-sw}\xspace}

% ---------------------------------------------------------------------------
% Code and identifiers.
% ---------------------------------------------------------------------------
% \code offers TeX no place to break, and this series is dense with identifiers
% long enough to overflow a line on their own (cds_list_for_each_entry_reverse,
% -DURCU_TXN_SW_EXCL_VALIDATE). Two of them in one sentence is enough to push
% text into the margin no matter how the surrounding prose is stretched.
%
% Since every such identifier already writes its underscores as \_, redefining
% \_ locally gives TeX a legal breakpoint after each one -- silently, with NO
% hyphen. That last part is the whole reason not to reach for hyphenat's [htt]
% option instead: a hyphen inside a C identifier is ambiguous, and the reader
% cannot tell whether it belongs to the name. For a paper whose purpose is
% enablement, an identifier someone has to guess at is a real defect; a slightly
% loose line is not.
%
% Scoped to \texttt's own group, so \_ keeps its normal meaning everywhere else.
%
% \let, not \renewcommand, and \DeclareRobustCommand, not \newcommand: \code is
% used inside \caption, which is a MOVING argument (written to the .aux and
% re-read). A fragile \code expands there and the \renewcommand form dies with
% "Argument of \caption@ydblarg has an extra }" -- a fatal error whose message
% names neither \code nor the caption it came from.
\newcommand{\codeunderscore}{\textunderscore\allowbreak}
\DeclareRobustCommand{\code}[1]{\texttt{\let\_\codeunderscore #1}}
\DeclareRobustCommand{\fn}[1]{\code{#1()}}

% A record in a transaction's write set: {slot, old, new}.
\newcommand{\record}[3]{$\{#1,\,#2 \rightarrow #3\}$}
% An edit on one slot, old to new.
\newcommand{\edit}[2]{$#1 \rightarrow #2$}

% Transaction status values.
\newcommand{\UNDECIDED}{\textsc{undecided}\xspace}
\newcommand{\SUCCEEDED}{\textsc{succeeded}\xspace}
\newcommand{\FAILED}{\textsc{failed}\xspace}
\newcommand{\ABORT}{\textsc{abort}\xspace}

% ---------------------------------------------------------------------------
% Editorial markers. Define \draftmode in the preamble to make these visible;
% they vanish in a submission build. Never ship with \draftmode set.
% ---------------------------------------------------------------------------
\newif\ifdraftmode
\draftmodefalse

% These use \color/\bfseries SWITCHES inside a group rather than
% \textcolor{}{}/\textbf{}, whose arguments are short: an editorial note
% spanning a blank line would otherwise die with "Paragraph ended before
% \@textcolor was complete". Notes routinely span paragraphs, so the switch
% form is the only one that survives real use.
\newcommand{\marker}[3]{\ifdraftmode{\color{#1}\bfseries[#2: #3]}\fi}

\newcommand{\TODO}[1]{\marker{red}{TODO}{#1}}
\newcommand{\NOTE}[1]{\marker{blue}{NOTE}{#1}}
% Claims needing a citation before submission.
\newcommand{\NEEDCITE}[1]{\marker{orange}{CITE}{#1}}
% Numbers currently sourced from source comments rather than a reproducible run.
\newcommand{\UNVERIFIED}[1]{\marker{purple}{UNVERIFIED}{#1}}
.
%
% Keep this class-agnostic: nothing here may assume `article`, so that a paper
% can be retargeted to acmart/IEEEtran by changing only its own preamble.

% ---------------------------------------------------------------------------
% Terminology.
%
% These exist so the series says the same thing the same way every time, and so
% that a terminology decision is changed in one place. The guarantee wording is
% load-bearing: this engine is NOT STM and is NOT serializable. See
% \pseudotxn below and the note in README.md.
% ---------------------------------------------------------------------------
\newcommand{\pseudotxn}{pseudo-transaction\xspace}
\newcommand{\pseudotxns}{pseudo-transactions\xspace}
\newcommand{\mcas}{MCAS\xspace}
\newcommand{\kcas}{$k$-CAS\xspace}
\newcommand{\rcu}{RCU\xspace}
\newcommand{\QSBR}{QSBR\xspace}
\newcommand{\urcu}{userspace RCU\xspace}
\newcommand{\sw}{single-writer\xspace}
\newcommand{\mw}{multi-writer\xspace}
\newcommand{\fliplatch}{flip-latch\xspace}
\newcommand{\ryw}{read-your-own-writes\xspace}

% Engine identifiers.
\newcommand{\urcutxn}{\texttt{urcu-txn}\xspace}
\newcommand{\rcumcas}{\texttt{rcu-mcas}\xspace}
\newcommand{\rcutxnsw}{\texttt{rcu-txn-sw}\xspace}

% ---------------------------------------------------------------------------
% Code and identifiers.
% ---------------------------------------------------------------------------
% \code offers TeX no place to break, and this series is dense with identifiers
% long enough to overflow a line on their own (cds_list_for_each_entry_reverse,
% -DURCU_TXN_SW_EXCL_VALIDATE). Two of them in one sentence is enough to push
% text into the margin no matter how the surrounding prose is stretched.
%
% Since every such identifier already writes its underscores as \_, redefining
% \_ locally gives TeX a legal breakpoint after each one -- silently, with NO
% hyphen. That last part is the whole reason not to reach for hyphenat's [htt]
% option instead: a hyphen inside a C identifier is ambiguous, and the reader
% cannot tell whether it belongs to the name. For a paper whose purpose is
% enablement, an identifier someone has to guess at is a real defect; a slightly
% loose line is not.
%
% Scoped to \texttt's own group, so \_ keeps its normal meaning everywhere else.
%
% \let, not \renewcommand, and \DeclareRobustCommand, not \newcommand: \code is
% used inside \caption, which is a MOVING argument (written to the .aux and
% re-read). A fragile \code expands there and the \renewcommand form dies with
% "Argument of \caption@ydblarg has an extra }" -- a fatal error whose message
% names neither \code nor the caption it came from.
\newcommand{\codeunderscore}{\textunderscore\allowbreak}
\DeclareRobustCommand{\code}[1]{\texttt{\let\_\codeunderscore #1}}
\DeclareRobustCommand{\fn}[1]{\code{#1()}}

% A record in a transaction's write set: {slot, old, new}.
\newcommand{\record}[3]{$\{#1,\,#2 \rightarrow #3\}$}
% An edit on one slot, old to new.
\newcommand{\edit}[2]{$#1 \rightarrow #2$}

% Transaction status values.
\newcommand{\UNDECIDED}{\textsc{undecided}\xspace}
\newcommand{\SUCCEEDED}{\textsc{succeeded}\xspace}
\newcommand{\FAILED}{\textsc{failed}\xspace}
\newcommand{\ABORT}{\textsc{abort}\xspace}

% ---------------------------------------------------------------------------
% Editorial markers. Define \draftmode in the preamble to make these visible;
% they vanish in a submission build. Never ship with \draftmode set.
% ---------------------------------------------------------------------------
\newif\ifdraftmode
\draftmodefalse

% These use \color/\bfseries SWITCHES inside a group rather than
% \textcolor{}{}/\textbf{}, whose arguments are short: an editorial note
% spanning a blank line would otherwise die with "Paragraph ended before
% \@textcolor was complete". Notes routinely span paragraphs, so the switch
% form is the only one that survives real use.
\newcommand{\marker}[3]{\ifdraftmode{\color{#1}\bfseries[#2: #3]}\fi}

\newcommand{\TODO}[1]{\marker{red}{TODO}{#1}}
\newcommand{\NOTE}[1]{\marker{blue}{NOTE}{#1}}
% Claims needing a citation before submission.
\newcommand{\NEEDCITE}[1]{\marker{orange}{CITE}{#1}}
% Numbers currently sourced from source comments rather than a reproducible run.
\newcommand{\UNVERIFIED}[1]{\marker{purple}{UNVERIFIED}{#1}}
.
%
% Keep this class-agnostic: nothing here may assume `article`, so that a paper
% can be retargeted to acmart/IEEEtran by changing only its own preamble.

% ---------------------------------------------------------------------------
% Terminology.
%
% These exist so the series says the same thing the same way every time, and so
% that a terminology decision is changed in one place. The guarantee wording is
% load-bearing: this engine is NOT STM and is NOT serializable. See
% \pseudotxn below and the note in README.md.
% ---------------------------------------------------------------------------
\newcommand{\pseudotxn}{pseudo-transaction\xspace}
\newcommand{\pseudotxns}{pseudo-transactions\xspace}
\newcommand{\mcas}{MCAS\xspace}
\newcommand{\kcas}{$k$-CAS\xspace}
\newcommand{\rcu}{RCU\xspace}
\newcommand{\QSBR}{QSBR\xspace}
\newcommand{\urcu}{userspace RCU\xspace}
\newcommand{\sw}{single-writer\xspace}
\newcommand{\mw}{multi-writer\xspace}
\newcommand{\fliplatch}{flip-latch\xspace}
\newcommand{\ryw}{read-your-own-writes\xspace}

% Engine identifiers.
\newcommand{\urcutxn}{\texttt{urcu-txn}\xspace}
\newcommand{\rcumcas}{\texttt{rcu-mcas}\xspace}
\newcommand{\rcutxnsw}{\texttt{rcu-txn-sw}\xspace}

% ---------------------------------------------------------------------------
% Code and identifiers.
% ---------------------------------------------------------------------------
% \code offers TeX no place to break, and this series is dense with identifiers
% long enough to overflow a line on their own (cds_list_for_each_entry_reverse,
% -DURCU_TXN_SW_EXCL_VALIDATE). Two of them in one sentence is enough to push
% text into the margin no matter how the surrounding prose is stretched.
%
% Since every such identifier already writes its underscores as \_, redefining
% \_ locally gives TeX a legal breakpoint after each one -- silently, with NO
% hyphen. That last part is the whole reason not to reach for hyphenat's [htt]
% option instead: a hyphen inside a C identifier is ambiguous, and the reader
% cannot tell whether it belongs to the name. For a paper whose purpose is
% enablement, an identifier someone has to guess at is a real defect; a slightly
% loose line is not.
%
% Scoped to \texttt's own group, so \_ keeps its normal meaning everywhere else.
%
% \let, not \renewcommand, and \DeclareRobustCommand, not \newcommand: \code is
% used inside \caption, which is a MOVING argument (written to the .aux and
% re-read). A fragile \code expands there and the \renewcommand form dies with
% "Argument of \caption@ydblarg has an extra }" -- a fatal error whose message
% names neither \code nor the caption it came from.
\newcommand{\codeunderscore}{\textunderscore\allowbreak}
\DeclareRobustCommand{\code}[1]{\texttt{\let\_\codeunderscore #1}}
\DeclareRobustCommand{\fn}[1]{\code{#1()}}

% A record in a transaction's write set: {slot, old, new}.
\newcommand{\record}[3]{$\{#1,\,#2 \rightarrow #3\}$}
% An edit on one slot, old to new.
\newcommand{\edit}[2]{$#1 \rightarrow #2$}

% Transaction status values.
\newcommand{\UNDECIDED}{\textsc{undecided}\xspace}
\newcommand{\SUCCEEDED}{\textsc{succeeded}\xspace}
\newcommand{\FAILED}{\textsc{failed}\xspace}
\newcommand{\ABORT}{\textsc{abort}\xspace}

% ---------------------------------------------------------------------------
% Editorial markers. Define \draftmode in the preamble to make these visible;
% they vanish in a submission build. Never ship with \draftmode set.
% ---------------------------------------------------------------------------
\newif\ifdraftmode
\draftmodefalse

% These use \color/\bfseries SWITCHES inside a group rather than
% \textcolor{}{}/\textbf{}, whose arguments are short: an editorial note
% spanning a blank line would otherwise die with "Paragraph ended before
% \@textcolor was complete". Notes routinely span paragraphs, so the switch
% form is the only one that survives real use.
\newcommand{\marker}[3]{\ifdraftmode{\color{#1}\bfseries[#2: #3]}\fi}

\newcommand{\TODO}[1]{\marker{red}{TODO}{#1}}
\newcommand{\NOTE}[1]{\marker{blue}{NOTE}{#1}}
% Claims needing a citation before submission.
\newcommand{\NEEDCITE}[1]{\marker{orange}{CITE}{#1}}
% Numbers currently sourced from source comments rather than a reproducible run.
\newcommand{\UNVERIFIED}[1]{\marker{purple}{UNVERIFIED}{#1}}
